\chapter{Apêndice A — Construção das Variáveis de Volatilidade}
\label{apendiceA}

Este apêndice apresenta a formulação das variáveis utilizadas ao longo da
dissertação: retornos logarítmicos, volatilidade realizada (RV), volatilidade
implícita (IV), prêmio de risco de volatilidade do Bitcoin (BVRP) e retornos
futuros. As expressões seguem práticas padrão em finanças quantitativas, com
adaptações relevantes ao mercado de criptoativos (negociação contínua 24/7).

\section{Retornos logarítmicos}

Dado o preço diário do Bitcoin $P_t$, define-se o retorno logarítmico como:
\[
r_t = \ln\left(\frac{P_t}{P_{t-1}}\right).
\]
O uso de retornos logarítmicos é conveniente por (i) permitir aditividade no
tempo e (ii) apresentar boa compatibilidade com modelos contínuos de preços.

Esses retornos são utilizados tanto na construção da volatilidade realizada
quanto na definição de retornos futuros em diferentes horizontes.

\section{Volatilidade Realizada (RV)}

A volatilidade realizada anualizada, construída a partir de retornos diários em
um horizonte de $h$ dias, é definida por:
\[
RV_{h,t} = \sqrt{\frac{365}{h}\sum_{i=1}^{h} r_{t-i+1}^{2}},
\]
em que o fator 365 anualiza a medida, considerando que o mercado de criptoativos
opera continuamente (24/7).

Ao longo da dissertação, utiliza-se majoritariamente $h=30$, de modo que:
\[
RV_{30,t} = \sqrt{\frac{365}{30}\sum_{i=1}^{30} r_{t-i+1}^{2}}.
\]

\subsection{Observação sobre alinhamento temporal}

A definição acima corresponde ao estimador \textit{backward-looking}, isto é,
calculado com retornos observados no passado (janela móvel encerrando em $t$).
Essa construção é adequada para análises empíricas e preditivas quando se deseja
utilizar apenas informação disponível até $t$. Em contextos de precificação, é
comum comparar IV em $t$ com uma RV \textit{ex-post} (realizada ao longo do
período futuro). Quando pertinente, a dissertação deixa explícito o alinhamento
adotado em cada exercício empírico.

\subsection{Estimadores alternativos de RV}

A literatura apresenta estimadores mais robustos para RV, sobretudo quando há
informação intradiária e/ou preocupação com ruído de microestrutura. Exemplos incluem:
\begin{itemize}
    \item \textbf{Parkinson} (baseado em máximos e mínimos intradiários);
    \item \textbf{Garman--Klass} (usa abertura, máximo, mínimo e fechamento);
    \item \textbf{Rogers--Satchell} (robusto a \textit{drift});
    \item \textbf{MedRV} (baseado em estatísticas medianas, resistente a ruído).
\end{itemize}

Apesar dessas alternativas, adota-se a versão clássica baseada em retornos diários
por sua simplicidade, transparência e consistência com a disponibilidade de dados
e com o horizonte de 30 dias utilizado na volatilidade implícita.

\section{Volatilidade Implícita (IV)}

A volatilidade implícita a 30 dias ($IV_{30,t}$) é obtida de duas formas:

\begin{enumerate}
    \item diretamente do índice DVOL da Deribit, utilizado como proxy da volatilidade
    implícita de curto prazo;
    \item quando necessário, via interpolação linear no tempo entre vencimentos adjacentes.
\end{enumerate}

Seja $T_1$ e $T_2$ (em dias) dois vencimentos observáveis tais que $T_1 < 30 < T_2$.
Define-se então:
\[
IV_{30,t} = w \cdot IV_{T_1,t} + (1-w)\cdot IV_{T_2,t},
\]
com peso:
\[
w = \frac{T_2 - 30}{T_2 - T_1}.
\]

Esse procedimento é análogo à interpolação de maturidade utilizada em índices de
volatilidade de mercado, como no caso do VIX.

\section{Prêmio de Risco de Volatilidade (BVRP)}

O Prêmio de Risco de Volatilidade do Bitcoin (BVRP), variável central da dissertação,
é definido como:
\[
BVRP_t = IV_{30,t} - RV_{30,t}.
\]

Em termos econômicos, $BVRP_t$ mede o diferencial entre a incerteza precificada
no mercado de opções (via IV) e a volatilidade efetivamente observada no mercado
à vista (via RV, construída em janela móvel). Valores positivos indicam que a
volatilidade implícita está acima da volatilidade realizada, consistente com a
interpretação de um prêmio associado à demanda por proteção. Valores negativos
podem ocorrer em episódios de choques, liquidações e movimentos abruptos em que
a volatilidade realizada supera a implícita.

\begin{itemize}
    \item \textbf{BVRP positivo:} IV $>$ RV, típico de mercados com demanda por proteção e/ou
    aversão ao risco elevada;
    \item \textbf{BVRP negativo:} RV $>$ IV, típico de episódios de surpresa (choques exógenos,
    liquidações em cascata, \textit{jumps});
    \item \textbf{Alta magnitude:} indica períodos de forte assimetria entre volatilidade
    precificada e observada, frequentemente associados a mudanças de regime.
\end{itemize}

\section{Retornos futuros}

Os retornos futuros utilizados para avaliação preditiva são definidos como retornos
logarítmicos acumulados no horizonte $h$:
\[
ret\_fut\_h(t) = \ln(P_{t+h}) - \ln(P_t).
\]

Em particular, são utilizados:
\[
ret\_fut\_1(t) = \ln(P_{t+1}) - \ln(P_t), \qquad
ret\_fut\_5(t) = \ln(P_{t+5}) - \ln(P_t), \qquad
ret\_fut\_20(t) = \ln(P_{t+20}) - \ln(P_t).
\]

Os horizontes escolhidos capturam:
\begin{itemize}
    \item \textbf{1 dia}: resposta de curtíssimo prazo e influência de ruído;
    \item \textbf{5 dias}: dinâmica semanal frequentemente analisada em criptoativos;
    \item \textbf{20 dias}: horizonte aproximadamente mensal, com maior suavização de ruído.
\end{itemize}

\section{Defasagem das variáveis e prevenção de \textit{look-ahead bias}}

Para assegurar validade empírica e evitar \textit{look-ahead bias}, a modelagem
preditiva utiliza apenas informação disponível até o instante de decisão. Em
particular, variáveis explicativas observadas no instante $t$ são, quando necessário,
defasadas de modo que:
\[
X_t = X^{\text{observado}}_{t-1},
\]
o que garante que:
\begin{itemize}
    \item preditores não utilizem informações do mesmo dia do alvo quando isso implicar
    vazamento temporal;
    \item previsões sejam realizadas apenas com informações disponíveis à época;
    \item a avaliação fora da amostra seja genuína e reprodutível.
\end{itemize}

\section{Resumo}

Este apêndice consolidou a formulação matemática das variáveis centrais da dissertação
--- retornos, $RV_{30,t}$, $IV_{30,t}$, $BVRP_t$ e retornos futuros ---
assegurando precisão e alinhamento com práticas usuais na literatura de volatilidade
e prêmio de risco. Essas definições fundamentam a construção dos \textit{datasets},
as análises econométricas e os exercícios preditivos discutidos nos capítulos principais.
